%%%%%%%%%%%%%%%%%%%%%%%%%%%%%%%%%%%%%%%%%%%%%%%%%%%%%%%%%%%%%%%%%%%%%%%%%
%
%	LaTeX File for PhD Thesis - Discussion
%
%%%%%%%%%%%%%%%%%%%%%%%%%%%%%%%%%%%%%%%%%%%%%%%%%%%%%%%%%%%%%%%%%%%%%%%%%

\chapter{Discussion}
\label{chap:discussion}

\section{Introduction}
This chapter provides a comprehensive discussion of the findings presented in this thesis. We interpret the results from the preceding chapters, analyze the implications of our work for the development of wireless foundation models, and honestly assess the limitations of our approach. The goal is to synthesize the individual contributions into a cohesive narrative and position our work within the broader context of wireless AI.

\section{Interpretation of Results}
The experimental results validate the effectiveness of our three-stage framework. This section delves into the reasons behind the success of each component.

\subsection{Why DenoMAE Works Well}
DenoMAE's success stems from its novel application of self-supervised learning principles to the wireless domain. The key factors contributing to its strong performance are:
\begin{itemize}
    \item \textbf{Multimodal Learning:} By jointly processing time-domain IQ samples and constellation diagrams, DenoMAE learns a richer, more robust representation than either modality alone could provide.
    \item \textbf{Noise-aware Representations:} Treating noise as an explicit modality to be modeled, rather than as a simple corruption to be removed, forces the model to learn the underlying structure of the signal in a more principled way. This leads to superior performance in low-SNR environments.
    \item \textbf{Constellation Diagrams as Effective Visual Representations:} Our work validates that constellation diagrams serve as an effective visual representation for wireless signals, enabling the successful application of powerful vision-based models like Vision Transformers.
\end{itemize}

\subsection{LLM Reasoning Advantages}
The Plug-and-Play LLM framework demonstrates a paradigm shift from traditional black-box classifiers to interpretable, white-box reasoning systems. Its advantages include:
\begin{itemize}
    \item \textbf{Training-Free Classification:} The ability to perform classification without any task-specific training or fine-tuning is a significant advantage, enabling rapid deployment and adaptation.
    \item \textbf{Interpretability and Explainability:} By leveraging the natural language reasoning capabilities of LLMs, the framework provides transparent and human-understandable classification decisions.
    \item \textbf{Open-Set Recognition:} The LLM's vast world knowledge provides a natural mechanism for handling out-of-distribution or unknown signal types, a critical capability for real-world applications.
\end{itemize}

\subsection{Meta-Learning Benefits}
The Meta-Task framework highlights the importance of adaptive, learnable regularization in few-shot learning. Its benefits include:
\begin{itemize}
    \item \textbf{Rapid Adaptation:} The framework enables models to adapt to new tasks with very few labeled examples, a crucial requirement for dynamic wireless environments.
    - \textbf{Method-Agnostic:} The Meta-Task framework is designed to be compatible with a wide range of meta-learning algorithms, making it a versatile tool for improving generalization.
    \item \textbf{Improved Generalization:} By learning to regularize on a per-task basis, the framework prevents overfitting and improves generalization to unseen tasks and conditions.
\end{itemize}

\subsection{Synergies Between Stages}
The true power of our framework lies in the synergy between its three stages:
\begin{enumerate}
    \item \textbf{DenoMAE} provides a powerful, general-purpose representation that serves as the foundation for the other stages.
    \item The \textbf{Plug-and-Play LLM} leverages these rich features to perform flexible, zero-shot classification and reasoning.
    \item \textbf{Meta-Task} enables the rapid adaptation of the learned representations to new, unseen tasks with minimal data.
\end{enumerate}
This modular, staged approach allows each component to specialize in its respective function while contributing to a cohesive and powerful overall system.

\section{Implications for Wireless Foundation Models}
This work represents a significant step toward the development of true foundation models for the wireless domain.

\subsection{Path Toward General-Purpose Wireless AI}
Our framework demonstrates that it is possible to build general-purpose models for wireless signals that can be adapted to a wide range of tasks without starting from scratch. This is a fundamental departure from the traditional approach of building specialized models for each specific task.

\subsection{Data Efficiency Gains}
A key achievement of this work is the demonstration of high performance with significantly less labeled data than traditional supervised methods. By leveraging self-supervised pre-training and meta-learning, we show that it is possible to build powerful models with as few as 10,000 labeled samples.

\subsection{Scalability Considerations}
While our experiments are conducted on 10-class datasets, the principles and architectures are designed to be scalable. The use of transformers and the modular nature of the framework provide a clear path toward scaling to larger datasets and a greater number of modulation classes.

\section{Comparison with State-of-the-Art}
Our methods, both individually and as an integrated framework, demonstrate competitive or superior performance compared to existing state-of-the-art approaches.

\subsection{Performance Benchmarking}
As shown in the results chapter, DenoMAE, PnP-Align, and Meta-Task all achieve significant performance gains over their respective baselines. The integrated framework, benefiting from the synergies between these components, is expected to set a new benchmark in the field.

\subsection{Novel Contributions}
The novelty of our work lies in the unique combination of self-supervised multimodal learning, LLM-based reasoning, and meta-learning for regularization. This holistic approach addresses the key challenges of representation, reasoning, and adaptation in a unified framework.

\subsection{Limitations of Existing Approaches}
Prior work in wireless signal classification has often focused on building highly specialized, black-box models that require large amounts of labeled data and do not generalize well to new tasks or environments. Our work overcomes these limitations by embracing the principles of foundation models.

\section{Limitations of This Work}
It is important to acknowledge the limitations of the current work to guide future research.

\subsection{Dataset Constraints}
Our experiments are primarily conducted on two 10-class datasets (DenoMAE 2.0 and a custom 10K dataset), with a third dataset for cross-domain validation. While these are sufficient to demonstrate the principles of our framework, further validation on larger and more diverse benchmarks is needed.

\subsection{Modulation Type Coverage}
The range of modulation types covered in our datasets is limited. Future work should aim to expand the pre-training and evaluation to include a much wider variety of modulation schemes to build a truly universal model.

\subsection{Real-World Deployment Challenges}
There is a significant gap between laboratory experiments and real-world deployment. Factors such as hardware impairments, complex channel effects, and the need for real-time processing are not fully addressed in this work.

\subsection{Computational Requirements}
While our methods are designed to be data-efficient, the pre-training of large models like DenoMAE and the use of large LLMs still require significant computational resources.

\section{Broader Impact}
The development of wireless foundation models has far-reaching implications for various fields.

\subsection{Cognitive Radio Applications}
Our work can enable more intelligent and autonomous cognitive radio systems that can sense, learn, and adapt to their spectral environment in real-time.

\subsection{Spectrum Monitoring}
Foundation models can be used for more effective spectrum monitoring, helping to identify unauthorized transmissions, detect interference, and ensure regulatory compliance.

\subsection{Military and Security Implications}
In the military and security domains, this technology can be applied to signal intelligence (SIGINT), electronic warfare (EW), and the development of more resilient and secure communication systems.

\section{Future Research Directions}
This thesis opens up several exciting avenues for future research.

\subsection{Scaling to More Modulation Types}
A natural next step is to scale the pre-training to include hundreds or even thousands of modulation types, moving closer to a universal signal classifier.

\subsection{Multi-Task Foundation Models}
Future work should explore extending the framework beyond classification to a wider range of tasks, such as signal separation, demodulation, and parameter estimation.


\subsection{Real-Time Implementation}
% TODO: Deployment considerations

\subsection{Transfer to Other Wireless Domains}
% TODO: Beyond AMC to other wireless problems

\section{Summary}
% TODO: Chapter summary
